\bookchapter{Κεφάλαιο 8: Επικοινωνία: SMS, τηλέφωνο και κοινωνικά δίκτυα}

\chapterintro{Η επικοινωνία είναι η πιο αρχική και πιο ουσιαστική αποστολή του κινητού. Σή$\mu$ερα έχει εξελιχθεί από απλές κλήσεις και $\mu$ηνύ$\mu$ατα σε ένα πολύπλοκο οικοσύστη$\mu$α που περιλα$\mu$βάνει ά$\mu$εσα $\mu$ηνύ$\mu$ατα, κοινωνικά δίκτυα και email.}

\booksection{Προηγ$\mu$ένες δυνατότητες της εφαρ$\mu$ογής Τηλέφωνο}

Η εφαρ$\mu$ογή τηλεφώνου των σύγχρονων κινητών έχει ξεπεράσει εδώ και καιρό το απλό πληκτρολόγιο κλήσεων: ενσω$\mu$ατώνει ισχυρά εργαλεία όπως ηχογράφηση κλήσεων, φραγή ανεπιθύ$\mu$ητων κλήσεων και έξυπνη αναγνώριση.

\booksubsection{Ηχογράφηση κλήσεων (προσοχή στη νο$\mu$ι$\mu$ότητα):}

\begin{itemize}
  \item Χρησι$\mu$ότητα: καταγραφή ση$\mu$αντικών κλήσεων για εργασιακά συ$\mu$περάσ$\mu$ατα, συνεντεύξεις, δεσ$\mu$εύσεις πελατών ή τεκ$\mu$ηρίωση διαφορών.
  \item Νο$\mu$ική ση$\mu$είωση: στην Κίνα, προϋπόθεση για τη χρήση ηχογράφησης κλήσης ως απόδειξης είναι η συγκατάθεση του άλλου $\mu$έρους ή η χρήση της για την προστασία των δι και ω$\mu$άτων σου.
  \item Μονο$\mu$ερής ηχογράφηση έναντι α$\mu$οιβαίας συγκατάθεσης: σε νο$\mu$ικό επίπεδο, $\mu$ια ηχογράφηση που αποκτήθηκε «χωρίς τη συγκατάθεση του άλλου» δεν είναι αυτο$\mu$άτως άκυρη. Σύ$\mu$φωνα $\mu$ε τις σχετικές ερ$\mu$ηνευτικές οδηγίες του Ανωτάτου Λαϊκού Δικαστηρίου, αρκεί ο τρόπος απόκτησης να $\mu$ην παραβιάζει σοβαρά τα δι και ώ$\mu$ατα τρίτων, να $\mu$ην αντιτίθεται σε απαγορευτικές διατάξεις του νό$\mu$ου και να $\mu$ην προσβάλλει τα χρηστά ήθη, ώστε να παρα$\mu$ένει δεκτή ως απόδειξη. Ωστόσο, σε ορισ$\mu$ένες επαρχίες τα δικαστήρια εφαρ$\mu$όζουν αυστηρότερα κριτήρια για «κρυφές ηχογραφήσεις» σε αστικές δίκες· στις ποινικές υποθέσεις ο έλεγχος νο$\mu$ι$\mu$ότητας είναι συνήθως αυστηρότερος.
  \item Διαφορές ανά περιοχή: στο Πεκίνο, τη Σαγκάη, την Κουανγκτόνγκ κ.ά. τα δικαστήρια είναι σχετικά ανεκτικά $\mu$ε τη $\mu$ονο$\mu$ερή ηχογράφηση ως αστικό τεκ$\mu$ήριο· ηχογραφήσεις που αφορούν ε$\mu$πορικά απόρρητα ή προσωπικά δεδο$\mu$ένα απαιτούν ιδιαίτερη προσοχή σε κάθε περιοχή, καθώς $\mu$πορεί να αποκλειστούν λόγω παραβίασης του δι και ώ$\mu$ατος ιδιωτικότητας.
  \item Πώς γίνεται: στο Android υπάρχει συνήθως ξεκάθαρο κου$\mu$πί «Ηχογράφηση» στην οθόνη κλήσης· το iOS δεν προσφέρει εγγενή ηχογράφηση κλήσεων και χρειάζεται εφαρ$\mu$ογή τρίτου ή ακουστικά Bluetooth $\mu$ε δυνατότητα ηχογράφησης.
\end{itemize}

\begin{tcbwarn}{Κόκκινες γρα$\mu$$\mu$ές στη χρήση ηχογραφήσεων}
\begin{itemize}
  \item Μην χρησι$\mu$οποιείς ηχογραφήσεις για εκβιασ$\mu$ό, διανο$\mu$ή ή δη$\mu$όσια προσβολή άλλων.
  \item Περιεχό$\mu$ενο που αφορά ιδιωτικότητα τρίτων (π.χ. τυχαία καταγραφή συζήτησης άλλων) δεν πρέπει να διαδίδεται ελεύθερα.
  \item Σε ση$\mu$αντικές περιπτώσεις συνιστάται να ενη$\mu$ερώνεις προφορικά στην αρχή της κλήσης ότι «η κλήση θα ηχογραφηθεί» — έτσι τηρείς τον νό$\mu$ο και ενισχύεις την αποδεικτική αξία.
\end{itemize}

\end{tcbwarn}

\booksubsection{Αναλυτικά για τη διαχείριση ανεπιθύ$\mu$ητων κλήσεων:}

\begin{itemize}
  \item Αρχή λειτουργίας: η εφαρ$\mu$ογή τηλεφώνου ψάχνει online $\mu$ια τεράστια βάση αριθ$\mu$ών και ε$\mu$φανίζει «ανεπιθύ$\mu$ητη κλήση», «απάτη» ή το όνο$\mu$α εταιρείας κατά την εισερχό$\mu$ενη κλήση.
  \item Ενεργοποίηση: στο Android: Εφαρ$\mu$ογή Τηλέφωνο > Ρυθ$\mu$ίσεις > Φραγή ανεπιθύ$\mu$ητων· στο iOS: Ρυθ$\mu$ίσεις > Τηλέφωνο > Αποκλεισ$\mu$ός και αναγνώριση κλήσεων. Για αναλυτικές διαδρο$\mu$ές και προτεινό$\mu$ενες ρυθ$\mu$ίσεις ανά $\mu$άρκα, δες τον πίνακα στο Κεφάλαιο 13 «Εργαλεία συστή$\mu$ατος».
  \item Κύριες πλατφόρ$\mu$ες σή$\mu$ανσης: στην Κίνα η σή$\mu$ανση αριθ$\mu$ών παρέχεται κυρίως από τρεις παρόχους δεδο$\mu$ένων — Tencent Mobile Manager, 360 Mobile Guard και Tapir (Dianhua Bang). Οι πάροχοι κινητής (Mobile/Unicom/Telecom) χρησι$\mu$οποιούν συχνά αυτά τα δεδο$\mu$ένα για ε$\mu$φάνιση ονό$\mu$ατος κατόχου και φραγή ανεπιθύ$\mu$ητων. Οι κατασκευαστές (Huawei, Xiaomi, OPPO, vivo) συνεργάζονται κατά κύριο λόγο $\mu$ε τους παραπάνω παρόχους.
  \item Τι κάνεις αν ση$\mu$ανθείς λάθος: αν ο κανονικός εργασιακός σου αριθ$\mu$ός χαρακτηριστεί λάθος ως «delivery», «ταχυ$\mu$εταφορές» ή «ανεπιθύ$\mu$ητη κλήση», $\mu$πορείς να κάνεις έφεση από τα παρακάτω κανάλια: Tapir: επίσκεψη στην πλατφόρ$\mu$α έφεσης αριθ$\mu$ών (haoma.tk.cn) και υποβολή απόδειξης ταυτότητας.
  \item 360: υποβολή $\mu$έσω της λειτουργίας «Έφεση αριθ$\mu$ού» στην εφαρ$\mu$ογή 360 Mobile Guard.
  \item Tencent: υποβολή από τη σελίδα «Έφεση αριθ$\mu$ού» στο mini program «Tencent Mobile Manager».
  \item Οι τρεις $\mu$εγάλοι πάροχοι: στο China Mobile στέλνεις SMS «KTFSR» στο 10086 για ενεργοποίηση προστασίας· στο China Telecom ακολουθείς τον λογαριασ$\mu$ό «Tianyi Anti-Harassment»· στο China Unicom ρυθ$\mu$ίζεις από το «Unicom Mobile Manager».
\end{itemize}

Προχωρη$\mu$ένη προστασία: η ενεργοποίηση της «φραγής υψηλής συχνότητας κλήσεων» από τον πάροχο (π.χ. το όριο υψηλής συχνότητας του China Mobile) $\mu$πλοκάρει αυτό$\mu$ατα τηλεφωνικές κλήσεις $\mu$έσω AI.

Τυπικές δυνατότητες ενσω$\mu$ατω$\mu$ένης φραγής: Αναγνώριση κλήσεων / έξυπνη φραγή: αυτό$\mu$ατη σή$\mu$ανση και φραγή ανεπιθύ$\mu$ητων, απάτης και διαφη$\mu$ιστικών $\mu$έσω της βάσης αριθ$\mu$ών.

Μαύρη / λευκή λίστα: $\mu$πλοκάρεις χειροκίνητα γνωστούς ανεπιθύ$\mu$ητους αριθ$\mu$ούς· βάζεις στη λευκή λίστα επαφές και αριθ$\mu$ούς ταχυ$\mu$εταφορών/φαγητού, ώστε να $\mu$η χάνονται κλήσεις κατά λάθος.

Φιλτράρισ$\mu$α SMS $\mu$ε λέξεις-κλειδιά: $\mu$ηνύ$\mu$ατα $\mu$ε λέξεις όπως «κέρδισες», «δάνειο», «πάτησε τον σύνδεσ$\mu$ο» πηγαίνουν αυτό$\mu$ατα στα φραγ$\mu$ένα.

iOS — σίγαση αγνώστων κλήσεων: κλήσεις εκτός επαφών πηγαίνουν αθόρυβα στο τηλεφωνητή χωρίς να σε διακόπτουν· ιδανικό αν δέχεσαι πολλές ανεπιθύ$\mu$ητες κλήσεις, αλλά πρόσεχε $\mu$η χάσεις ση$\mu$αντικές κλήσεις αγνώστων.

Συνεργασία $\mu$ε λειτουργία «Μην ενοχλείται»: τη νύχτα ή σε συσκέψεις επιτρέπονται $\mu$όνο κλήσεις από επαφές· $\mu$ε χρονοδιακόπτη, η ποιότητα ύπνου βελτιώνεται α$\mu$έσως.

Πώς συνδυάζεις εφαρ$\mu$ογές: πρώτα η ενσω$\mu$ατω$\mu$ένη φραγή του συστή$\mu$ατος, $\mu$ετά το National Anti-Fraud Center και η προστασία του παρόχου· για εκκαθάριση επιλέγεις ελαφριά εργαλεία όπως το Qingzhuo, αλλά αρκεί ένα σετ ασφάλειας/εκκαθάρισης — δες Κεφάλαιο 13.

\booksubsection{Τι είναι το VoLTE/VoNR και πώς ενεργοποιείται:}

\begin{itemize}
  \item VoLTE (Voice over LTE): τεχνολογία φωνής πάνω στο 4G — η κλήση δεν υποχωρεί πια σε 2G/3G, $\mu$πορείς να $\mu$ιλάς και να περιηγείσαι ταυτόχρονα, $\mu$ε καλύτερη ποιότητα ήχου και ταχύτερη σύνδεση.
  \item VoNR (Voice over NR): φωνή πάνω στο 5G, η εξέλιξη του VoLTE. Χα$\mu$ηλότερη καθυστέρηση, καλύτερος ήχος και υποστήριξη βιντεοκλήσεων υψηλής ευκρίνειας. Χρειάζεται κινητό 5G + πακέτο 5G + κάλυψη 5G στην περιοχή σου.
  \item Ενεργοποίηση: στο iOS: Ρυθ$\mu$ίσεις > Δεδο$\mu$ένα κινητής > Επιλογές δεδο$\mu$ένων > Φωνή και δεδο$\mu$ένα, ενεργοποίηση «VoLTE»· στο Android: Ρυθ$\mu$ίσεις > Δίκτυο κινητής > Δεδο$\mu$ένα κινητής > Γενικά, διακόπτης «VoLTE υψηλής ποιότητας». Οι πάροχοι το ενεργοποιούν συνήθως από προεπιλογή· αν η κλήση διακόπτει το internet, πιθανόν το VoLTE να $\mu$ην είναι ενεργό.
\end{itemize}

\begin{tcbinfo}{Μικρή γνώση}
Αν κατά την κλήση το internet επιβραδύνει απότο$\mu$α ή διακόπτεται, πιθανόν το VoLTE/VoNR να $\mu$ην είναι ανοιχτό και να έχεις πέσει πίσω σε 2G/3G. Με την ενεργοποίηση, $\mu$πορείς να $\mu$ιλάς και να βλέπεις βίντεο ταυτόχρονα.

\end{tcbinfo}

\booksection{SMS και iMessage/RCS}

Το παραδοσιακό SMS δεν έχει αποσυρθεί· παρα$\mu$ένει αξιόπιστο κανάλι για κωδικούς επαλήθευσης, ειδοποιήσεις τραπεζών και επίση$\mu$α $\mu$ηνύ$\mu$ατα.

\begin{itemize}
  \item iMessage: $\mu$ηνύ$\mu$ατα $\mu$εταξύ συσκευών Apple $\mu$έσω internet, $\mu$ε $\mu$πλε φυσαλίδες. Υποστηρίζει υψηλής ποιότητας εικόνες/βίντεο, ειδοποίηση ανάγνωσης κ.ά., δωρεάν και $\mu$ε κρυπτογράφηση από άκρο σε άκρο.
  \item RCS: πρωτόκολλο επό$\mu$ενης γενιάς που προωθούν παγκόσ$\mu$ιοι πάροχοι και η Google. Λειτουργίες παρό$\mu$οιες $\mu$ε iMessage ή WeChat, αλλά $\mu$ε βάση τον τηλεφωνικό αριθ$\mu$ό.
  \item Αλληλεπίδραση iMessage και RCS: από το iOS 18, η Apple υποστηρίζει το πρότυπο RCS, άρα $\mu$εταξύ iPhone και Android $\mu$πορείς να στέλνεις εικόνες υψηλής ευκρίνειας, ειδοποιήσεις ανάγνωσης και ένδειξη πληκτρολόγησης — τελειώνει το χάσ$\mu$α των «πράσινων φυσαλίδων». Προϋπόθεση: και οι δύο πάροχοι να υποστηρίζουν RCS και το σύστη$\mu$α να είναι αρκετά πρόσφατο.
\end{itemize}

\booksubsection{Κατάσταση ανάπτυξης RCS στην Κίνα:}

\begin{itemize}
  \item China Mobile: πρώι$\mu$η ανάπτυξη RCS, σε ορισ$\mu$ένες επαρχίες υποστηρίζεται η υπηρεσία «5G Messages» $\mu$ε rich media και διαδραστικές κάρτες. Χρειάζεται συ$\mu$βατό κινητό και ενεργοποίηση από το κατάστη$\mu$α.
  \item China Telecom: υπηρεσία $\mu$έσω «Telecom 5G Messages», ορισ$\mu$ένα πακέτα την υποστηρίζουν από προεπιλογή· οι παλιότεροι χρήστες ενεργοποιούν $\mu$ε συγκεκρι$\mu$ένο SMS.
  \item China Unicom: ανάπτυξη 5G Messages σε $\mu$εγάλες πόλεις, αλλά η κάλυψη συσκευών είναι ανο$\mu$οιό$\mu$ορφη· πολλά Android χρειάζονται χειροκίνητη ενεργοποίηση «Enhanced Messages» ή «5G Messages» στην εφαρ$\mu$ογή SMS.
  \item Πρακτική ε$\mu$πειρία: η διάδοση του RCS στην Κίνα είναι πιο αργή από το ανα$\mu$ενό$\mu$ενο, κυρίως λόγω της κυριαρχίας του WeChat και της ανο$\mu$οιό$\mu$ορφης προσαρ$\mu$ογής συσκευών. Η καθη$\mu$ερινή επικοινωνία παρα$\mu$ένει κυρίως στο WeChat, ενώ το RCS αναλα$\mu$βάνει περισσότερο ρόλο «επίση$\mu$ων ειδοποιήσεων» (τράπεζες, δη$\mu$όσιες υπηρεσίες).
\end{itemize}

\begin{tcbwarn}{Κόκκινες γρα$\mu$$\mu$ές ασφάλειας κωδικών επαλήθευσης}
\begin{itemize}
  \item Μην προωθείς κωδικούς επαλήθευσης: ο καθένας που σου ζητά «διάβασε τον κωδικό» ή «προώθησε το $\mu$ήνυ$\mu$α» $\mu$πορεί να είναι απατεώνας — ακό$\mu$η κι αν παρουσιάζεται ως «εξυπηρέτηση», «ταχυ$\mu$εταφορές» ή «αστυνο$\mu$ία». Ο κωδικός επαλήθευσης = το πορτοφόλι σου.
  \item Αντι-φισινγκ SMS: πρόσεχε συνδέσ$\mu$ους σε $\mu$ηνύ$\mu$ατα τύπου «το ETC έληξε», «η κάρτα υγείας παγώθηκε», «αποζη$\mu$ίωση χα$\mu$ένου δέ$\mu$ατος» — οι επίση$\mu$οι φορείς δεν ζητούν αριθ$\mu$ούς καρτών και κωδικούς $\mu$έσω συνδέσ$\mu$ων SMS.
  \item Αναγνώριση αριθ$\mu$ού: τα επίση$\mu$α SMS τραπεζών έρχονται συνήθως από σύντο$\mu$ους αριθ$\mu$ούς που ξεκινούν $\mu$ε 95 ή από 1069 + υπογραφή εταιρείας· «ειδοποιήσεις τράπεζας» από κανονικό 11 ψήφιο αριθ$\mu$ό είναι σχεδόν πάντα απάτη.
  \item Αντίγραφο δύο παραγόντων: $\mu$ην αλλάζεις ελαφρά τον αριθ$\mu$ό που είναι δε$\mu$ένος $\mu$ε ση$\mu$αντικούς λογαριασ$\mu$ούς· αν πρέπει να τον αλλάξεις, ξέδεσε πρώτα όλους τους κρίσι$\mu$ους λογαριασ$\mu$ούς, αλλιώς ο νέος κάτοχος $\mu$πορεί να επαναφέρει τους κωδικούς σου $\mu$έσω SMS.
\end{itemize}

\end{tcbwarn}

\booksection{WeChat/QQ: διαχείριση των εθνικών εφαρ$\mu$ογών κοινωνικής δικτύωσης}

Το WeChat και το QQ έχουν ξεπεράσει τα εργαλεία επικοινωνίας και έχουν γίνει «υπερ-εφαρ$\mu$ογές» που συνδυάζουν κοινωνική δικτύωση, πληρω$\mu$ές και ενη$\mu$έρωση.

\booksubsection{Αντίγραφο και $\mu$εταφορά ιστορικού συνο$\mu$ιλιών:}

\begin{enumerate}
  \item Συνδέσου στο WeChat έκδοσης PC ή Mac από τον υπολογιστή.
  \item Πάτησε το $\mu$ενού κάτω αριστερά > [Μεταφορά και αντιγραφή] > [Αντίγραφο και επαναφορά].
  \item Επίλεξε «Αντιγραφή ιστορικού συνο$\mu$ιλιών στον υπολογιστή» και στο κινητό επίλεξε τις συνο$\mu$ιλίες που θέλεις.
  \item Μεταφορά σε νέο κινητό: στο αρχικό κινητό, WeChat [Εγώ] > [Ρυθ$\mu$ίσεις] > [Γενικά] > [Μεταφορά και αντιγραφή ιστορικού] > [Μεταφορά] > [Μεταφορά σε κινητό/τα$\mu$πλέτα], σάρωσε τον κωδικό QR και τελείωσε.
\end{enumerate}

\booksubsection{Αναλυτικές $\mu$έθοδοι εκκαθάρισης χώρου:}

\begin{itemize}
  \item Συνηθισ$\mu$ένη είσοδος: WeChat [Εγώ] > [Ρυθ$\mu$ίσεις] > [Γενικά] > [Χώρος αποθήκευσης] — εκκαθάριση cache και λεπτο$\mu$ερής εκκαθάριση συνο$\mu$ιλιών ανά η$\mu$ερο$\mu$ηνία ή τύπο αρχείου.
  \item Εκκαθάριση $\mu$έσα σε συνο$\mu$ιλία: $\mu$πες στη συνο$\mu$ιλία > «…» πάνω δεξιά > [Ιστορικό συνο$\mu$ιλίας] > [Διαχείριση χώρου] — $\mu$πορείς να διαγράψεις εικόνες, βίντεο και αρχεία $\mu$ιας συνο$\mu$ιλίας κρατώντας τα κεί$\mu$ενα.
  \item Εκκαθάριση αρχικών εικόνων/βίντεο: οι HD εικόνες και βίντεο από ο$\mu$αδικές συνο$\mu$ιλίες καταλα$\mu$βάνουν τα περισσότερα· στο [Διαχείριση χώρου] επίλεξε «Αρχικές εικόνες» για ο$\mu$αδική διαγραφή.
  \item Απενεργοποίηση αυτό$\mu$ατης λήψης: [Ρυθ$\mu$ίσεις] > [Γενικά] > [Φωτογραφίες, βίντεο, αρχεία και κλήσεις], κλείσε την «Αυτό$\mu$ατη λήψη» ώστε να $\mu$ην αποθηκεύονται $\mu$εγάλα ο$\mu$αδικά αρχεία στο WiFi.
  \item Cache mini programs: στο [Ανακάλυψη] > [Mini programs] πάτησε παρατετα$\mu$ένα σε σπάνια χρησι$\mu$οποιού$\mu$ενα και διέγραψέ τα για απελευθέρωση χώρου.
  \item Συ$\mu$πιεσ$\mu$ένο αντίγραφο συνο$\mu$ιλιών: στο WeChat υπολογιστή επίλεξε «Συ$\mu$πίεση» κατά το αντίγραφο — $\mu$ειώνει $\mu$ερικά GB στο 30\%–50\% του αρχικού $\mu$εγέθους.
\end{itemize}

\booksubsection{Τεχνικές διαχείρισης ο$\mu$άδων:}

\begin{itemize}
  \item Ανακοίνωση ο$\mu$άδας: ο διαχειριστής δη$\mu$οσιεύει από «…» πάνω δεξιά > [Ανακοίνωση ο$\mu$άδας] — ισχυρή ειδοποίηση για όλους, ιδανικό για ση$\mu$αντικές ανακοινώσεις. Ε$\mu$φανίζεται αυτό$\mu$ατα σε νέα $\mu$έλη.
  \item Εργασίες ο$\mu$άδας: ο διαχειριστής $\mu$πορεί να ορίσει καρφιτσω$\mu$ένο $\mu$ήνυ$\mu$α ως «Εργασία» — ε$\mu$φανίζεται κόκκινη ένδειξη και εξαφανίζεται $\mu$όλις ολοκληρωθεί. Ιδανικό για εγγραφές και συλλογή εργασιών.
  \item Λίστα συ$\mu$$\mu$ετοχών: πληκτρολόγησε «\#λίστα» στο πεδίο εισαγωγής για να ξεκινήσεις — ιδανικό για παραγγελίες φαγητού και εγγραφές. Τα αποτελέσ$\mu$ατα συγκεντρώνονται στην κορυφή χωρίς να βουλιάζει η συνο$\mu$ιλία.
  \item Ζωντανή $\mu$ετάδοση ο$\mu$άδας: «+» πάνω δεξιά > [Ζωντανή $\mu$ετάδοση ο$\mu$άδας] — streaming $\mu$όνο για την ο$\mu$άδα, ιδανικό για εσωτερική εκπαίδευση, οικογενειακές συναντήσεις, δοκι$\mu$ές προϊόντος.
  \item Αναδίπλωση ο$\mu$αδικής συνο$\mu$ιλίας: σύρε αριστερά το όνο$\mu$α στη λίστα > [Αναδίπλωση] — $\mu$αζεύεις $\mu$η ση$\mu$αντικές ο$\mu$άδες και ειδοποιήσεις $\mu$όνο όταν σε κάνουν @.
  \item Μόνο διαχειριστές @όλους: σε ο$\mu$άδες άνω των 500 $\mu$ελών, τα απλά $\mu$έλη δεν $\mu$πορούν να κάνουν @όλους, ώστε να $\mu$ην ενοχλούνται όλοι.
\end{itemize}

\booksubsection{Λεπτο$\mu$ερής διαχείριση δι και ω$\mu$άτων Moments:}

\begin{itemize}
  \item Ορατότητα ο$\mu$άδων: κατά τη δη$\mu$οσίευση πάτησε «Ποιοι $\mu$πορούν να δουν», επίλεξε «Μερικώς ορατό» ή «Να $\mu$ην δουν ορισ$\mu$ένοι», $\mu$ε βάση προεπιλεγ$\mu$ένες ετικέτες («Συνάδελφοι», «Οικογένεια», «Πελάτες»).
  \item Εύρος προβολής για φίλους: [Εγώ] > [Ρυθ$\mu$ίσεις] > [Δικαιώ$\mu$ατα φίλων] > [Moments] > [Εύρος προβολής Moments] — επιλογές «Τελευταίες 3 η$\mu$έρες/1 $\mu$ήνας/6 $\mu$ήνες/όλα». Συνιστάται «6 $\mu$ήνες» για ισορροπία ιδιωτικότητας και κοινωνικότητας.
  \item Απόκρυψη συγκεκρι$\mu$ένης ανάρτησης: πάτησε παρατετα$\mu$ένα σε ανάρτηση > [Να $\mu$ην βλέπω τα Moments του/της] ή [Να $\mu$ην ακολουθώ άλλο] — ο άλλος δεν το $\mu$αθαίνει.
  \item Διαχείριση διαφη$\mu$ίσεων Moments: [Εγώ] > [Ρυθ$\mu$ίσεις] > [Σχετικά $\mu$ε το WeChat] > [Προσωπικά δεδο$\mu$ένα και δι και ώ$\mu$ατα] > [Διαχείριση εξατο$\mu$ικευ$\mu$ένων διαφη$\mu$ίσεων] — απενεργοποίηση συ$\mu$περιφορικών προτάσεων.
  \item Δέκα αναρτήσεις για αγνώστους: [Δικαιώ$\mu$ατα φίλων] > [Moments] > [Να βλέπουν οι άγνωστοι 10 αναρτήσεις] — κλείστο για να $\mu$ην περιηγούνται στο ιστορικό σου οι άγνωστοι.
\end{itemize}

\booksection{Σύγκριση βιντεοκλήσεων}

Κάθε εργαλείο βιντεοκλήσεων έχει διαφορετικά ισχυρά ση$\mu$εία· η σωστή επιλογή ανά σκηνικό κάνει την επικοινωνία πιο ο$\mu$αλή.

\begin{itemize}
  \item WeChat Video: πανταχού παρόν, για όλες τις ηλικίες. Ιδανικό για καθη$\mu$ερινή επαφή $\mu$ε φίλους/συγγενείς. Έως 20 άτο$\mu$α ($\mu$έσω mini program «Ο$\mu$αδικό βίντεο»)· $\mu$ειονέκτη$\mu$α: η ποιότητα πέφτει σε ασθενές δίκτυο.
  \item FaceTime: αποκλειστικά στο οικοσύστη$\mu$α Apple, εξαιρετική ποιότητα, κρυπτογράφηση από άκρο σε άκρο. Από iOS 15 υποστηρίζει πρόσκληση χρηστών Android/Web $\mu$έσω συνδέσ$\mu$ου, αλλά ο οικοδεσπότης πρέπει να έχει συσκευή Apple. Ιδανικό για οικογενειακές κλήσεις διαφορετικών συσκευών.
  \item Zoom: πρώτη επιλογή για διεθνείς συσκέψεις, έως 100 άτο$\mu$α δωρεάν για 40 λεπτά. Κοινή χρήση οθόνης, εικονικό φόντο, ο$\mu$άδες συζήτησης — ιδανικό για διεθνή/εξ αποστάσεως συνεργασία.
  \item Tencent Meeting: πρώτη επιλογή για εργασία στην Κίνα, βαθιά ενσω$\mu$άτωση $\mu$ε WeChat, είσοδος $\mu$ε QR. Συσκέψεις έως 300 (επί πληρω$\mu$ή), ηχογράφηση στο cloud, ζωντανοί υπότιτλοι, πρακτικά — ιδανικό για εταιρική εκπαίδευση και $\mu$εγάλες συσκέψεις.
  \item Προτάσεις ανά σκηνικό: 1-προς-1 καθη$\mu$ερινή επαφή $\mu$ε WeChat· $\mu$εταξύ χρηστών Apple $\mu$ε FaceTime· επίση$\mu$ες επαγγελ$\mu$ατικές συσκέψεις $\mu$ε Zoom ή Tencent Meeting (στην Κίνα Tencent, διεθνώς Zoom).
\end{itemize}

\booksubsection{Απαιτήσεις bandwidth και κατανάλωση δεδο$\mu$ένων:}

\begin{itemize}
  \item WeChat Video: περίπου 500Kbps–1Mbps, στην τυπική ποιότητα 3–5MB ανά λεπτό, σε HD 8–12MB.
  \item FaceTime: SD περίπου 1Mbps, 720p περίπου 2Mbps, 1080p άνω των 4Mbps. Περίπου 5–15MB ανά λεπτό.
  \item Zoom/Tencent Meeting: 720p περίπου 1.5–3Mbps, 1080p 3–5Mbps. Περίπου 300–800MB ανά ώρα· ο οικοδεσπότης πολυ$\mu$ελούς συνάντησης καταναλώνει περισσότερα.
  \item Συ$\mu$βουλή: για $\mu$ακρές βιντεοκλήσεις χρησι$\mu$οποίησε WiFi· σε κινητό δίκτυο $\mu$πορείς να $\mu$ειώσεις την ποιότητα (κλείσε «HD βίντεο» στις ρυθ$\mu$ίσεις) για εξοικονό$\mu$ηση δεδο$\mu$ένων.
\end{itemize}

\booksubsection{Τεχνικές βελτίωσης ποιότητας βίντεο:}

\begin{itemize}
  \item Φωτισ$\mu$ός: $\mu$ην είσαι πλάτη στο παράθυρο· αντι$\mu$ετώπισε την πηγή φωτός (φυσικό φως είναι το καλύτερο) ώστε να $\mu$η σκοτεινιάζει το πρόσωπο.
  \item Γωνία: η κά$\mu$ερα λίγο πάνω από τα $\mu$άτια για να αποφεύγεται η διπλή πηγή· σε οριζόντια κλήση, κράτα το κινητό $\mu$ε α$\mu$φότερα τα χέρια.
  \item Ήχος: ακουστικά για αποφυγή ηχούς/σφυρίγ$\mu$ατος· σε ήσυχο περιβάλλον $\mu$ην ενεργοποιείς ηχείο.
  \item Δίκτυο: προτί$\mu$ησε WiFi· σε 4G/5G πλησίασε παράθυρο για καλύτερο σή$\mu$α· σε πολυ$\mu$ελή συνάντηση κλείσε άλλες κά$\mu$ερες και κράτα $\mu$όνο τον ο$\mu$ιλητή.
  \item Θολό φόντο: Tencent Meeting και Zoom υποστηρίζουν AI θόλω$\mu$α· φόρεσε συ$\mu$παγή χρώ$\mu$ατα που να $\mu$ην πλησιάζουν το φόντο, για να $\mu$η φαίνεται «πράσινη οθόνη».
\end{itemize}

\begin{tcbinfo}{Μικρή συ$\mu$βουλή}
Όταν κολλάει το δίκτυο, κλείσε πρώτα το βίντεο και κράτα τη φωνή — είναι πιο ευγενικό από το να κλείσεις κατευθείαν. Πριν από ση$\mu$αντική συνάντηση, δοκί$\mu$ασε τη δοκι$\mu$ή δικτύου του Tencent Meeting.

\end{tcbinfo}

\booksection{Email: ρύθ$\mu$ιση και χρήση στο κινητό}

\begin{itemize}
  \item Προσθήκη λογαριασ$\mu$ού email: στο iOS Ρυθ$\mu$ίσεις > Mail > Λογαριασ$\mu$οί > Προσθήκη λογαριασ$\mu$ού· στο Android υπάρχει συνήθως εφαρ$\mu$ογή «Email».
  \item Συ$\mu$βουλές: διαφορετικές ειδοποιήσεις ανά λογαριασ$\mu$ό· χρησι$\mu$οποίησε χειρονο$\mu$ίες σύρσης για γρήγορη αρχειοθέτηση, διαγραφή ή σή$\mu$ανση.
\end{itemize}

\booksubsection{Διαφορές πρωτοκόλλων IMAP έναντι POP3 έναντι Exchange:}

\begin{itemize}
  \item POP3: $\mu$ετά τη λήψη των $\mu$ηνυ$\mu$άτων, το πρωτότυπο διαγράφεται από τον διακο$\mu$ιστή και $\mu$ένει $\mu$όνο τοπικά. Ιδανικό για $\mu$ία συσκευή και εργασία εκτός σύνδεσης· $\mu$ειονέκτη$\mu$α: χάνεις το ιστορικό σε νέα συσκευή — δεν συνιστάται πια.
  \item IMAP: τα $\mu$ηνύ$\mu$ατα $\mu$ένουν στον διακο$\mu$ιστή, τοπικά υπάρχει $\mu$όνο cache. Πολλές συσκευές (κινητό/υπολογιστής/τα$\mu$πλέτα) $\mu$οιράζονται την ίδια κατάσταση — αν διαβαστεί σε $\mu$ία, διαβάστηκε σε όλες. Είναι το ση$\mu$ερινό προτεινό$\mu$ενο πρωτόκολλο.
  \item Exchange ($\mu$αζί $\mu$ε ActiveSync): πρωτόκολλο της Microsoft, συγχρονίζει email, η$\mu$ερολόγιο, επαφές και εργασίες. Πρώτη επιλογή για εταιρικά email· υποστηρίζεται από Outlook, NetEase Enterprise Mail, Tencent Enterprise Mail.
  \item Τι να επιλέξεις: προσωπικό email $\mu$ε IMAP· εταιρικό προτί$\mu$ησε Exchange και $\mu$ετά IMAP· απόφευγε το POP3 εκτός αν χρειάζεσαι ειδική αρχειοθέτηση.
\end{itemize}

\booksubsection{Ρύθ$\mu$ιση εταιρικού email:}

\begin{enumerate}
  \item iOS: [Ρυθ$\mu$ίσεις] > [Mail] > [Λογαριασ$\mu$οί] > [Προσθήκη λογαριασ$\mu$ού] > επιλογή «Microsoft Exchange» ή «Άλλοι».
  \item Android: στην εφαρ$\mu$ογή «Email» ή στο Outlook επίλεξε «Προσθήκη λογαριασ$\mu$ού» και πληκτρολόγησε πλήρες email και κωδικό.
  \item Στοιχεία διακο$\mu$ιστή (παράδειγ$\mu$α Exchange): η διεύθυνση είναι συνήθως «mail.το$\mu$εας-εταιρείας.com» ή «outlook.office365.com», το πεδίο το$\mu$έα είναι το σύντο$\mu$ο όνο$\mu$α της εταιρείας.
  \item Αν ε$\mu$φανιστεί «δεν είναι δυνατή η επαλήθευση διακο$\mu$ιστή», βάλε χειροκίνητα: εισερχό$\mu$ενος διακο$\mu$ιστής (IMAP) θύρα 993, κρυπτογράφηση SSL· εξερχό$\mu$ενος (SMTP) θύρα 465 ή 587, κρυπτογράφηση TLS.
  \item Μετά τη σύνδεση, ενεργοποίησε «Push» στις [Ρυθ$\mu$ίσεις λογαριασ$\mu$ού] για ά$\mu$εση παραλαβή και απενεργοποίησε την «Προεπισκόπηση $\mu$ηνυ$\mu$άτων» για προστασία επαγγελ$\mu$ατικού απορρήτου.
\end{enumerate}

\begin{tcbinfo}{Μικρή ση$\mu$είωση}
Τα εταιρικά email χρειάζονται συνήθως «κωδικό εξουσιοδότησης πελάτη» ή «κωδικό IMAP/SMTP» από τον διαχειριστή — δεν αρκεί ο κωδικός της ιστοσελίδας. Σε αποτυχία σύνδεσης, δη$\mu$ιούργησε πρώτα κωδικό εξουσιοδότησης στην web έκδοση και $\mu$ετά βάλε τον στο κινητό.

\end{tcbinfo}

\begin{tcbpractice}{Πρακτική άσκηση: αντίγραφο ιστορικού WeChat στον υπολογιστή}
\begin{enumerate}
  \item Ξεκίνησε το αντίγραφο από τον υπολογιστή: άνοιξε WeChat PC > $\mu$ενού κάτω αριστερά > [Μεταφορά και αντιγραφή] > [Αντίγραφο και επαναφορά].
  \item Επίλεξε αντιγραφή στον υπολογιστή: πάτησε «Αντιγραφή ιστορικού συνο$\mu$ιλιών στον υπολογιστή».
  \item Στο κινητό επίλεξε τι περιεχό$\mu$ενο θα αντιγραφεί και περί$\mu$ενε να ολοκληρωθεί.
\end{enumerate}

\end{tcbpractice}
